\documentclass[11pt]{article}
\usepackage{amsmath,amssymb}
\usepackage[margin=1in]{geometry}
\begin{document}

\subsection*{Equal-combination response as a function of $p$}

Model (paper, restricted to the two directions as the only features, weights
$w_1,w_2$, in-plane angle $\varphi$, and $g \equiv \sin^p\!\alpha\,\|a\|^p$):
\[
R(\varphi) = F\!\big(\, S(\varphi) \,\big),
\qquad
S(\varphi) = g\,\big[\, w_1\cos^p\varphi + w_2\sin^p\varphi \,\big].
\]

The three measured conditions (at fixed perturbation angle $\alpha$):
\[
\underbrace{S_1 = g\,w_1}_{\varphi=0},
\qquad
\underbrace{S_2 = g\,w_2}_{\varphi=\pi/2},
\qquad
\underbrace{S_{\mathrm{mix}} = g\,(w_1+w_2)\,2^{-p/2}}_{\varphi=\pi/4,\ \cos\varphi=\sin\varphi=2^{-1/2}}.
\]

Hence, independent of $w_1,w_2,g$:
\[
\boxed{\ \frac{S_{\mathrm{mix}}}{S_1+S_2} = 2^{-p/2}\ }
\qquad\text{and, for }w_1=w_2:\quad
\frac{S_{\mathrm{mix}}}{S_{\mathrm{axis}}} = 2^{\,1-p/2}.
\]

So $2^{-p/2}\in(0,\tfrac12]$ for $p\ge 2$, and $\to 0$ only as $p\to\infty$.
Identifying $R\approx S$ ($F$ locally linear), the proxy
$R_{\mathrm{mix}}/(R_1+R_2)$ \emph{equals} $2^{-p/2}$: minimizing it
$\Leftrightarrow$ pushing $p\to\infty$.

\subsection*{Why $R_{\mathrm{mix}}\to 0$ is not the $p\to\infty$ box}

The relation above fixes only the \emph{single} diagonal point at the axis
magnitude. The full $p\to\infty$ contour (the box) still \emph{breaks} along the
diagonal --- at the corner. In normalized coordinates
$(x,y)=(\sin\alpha\cos\varphi/\sin\theta_1,\ \sin\alpha\sin\varphi/\sin\theta_2)$
the superellipse is $|x|^p+|y|^p=1$; the diagonal ($x=y$) crosses it at
\[
x = y = 2^{-1/p},
\qquad\text{so as }p\to\infty:\quad x=y\to 1,
\]
i.e.\ the corner $(\theta_1,\theta_2)$, at geodesic radius $\sqrt{2}\,\theta$
(equal intercepts). So even for $p=\infty$ the diagonal breaks at a
\emph{finite} angle $\sqrt2\,\theta$; for any finite $p$ the diagonal-point
ratio $2^{-p/2}$ is strictly positive.

The optimizer instead drives $R_{\mathrm{mix}}\approx 0$ \emph{sustained} far
past $\sqrt2\,\theta$ (data along the diagonal out to $\gtrsim 60^\circ$ with
$\theta\approx 22^\circ$, corner $\approx 31^\circ$). No $p\in(0,\infty]$
produces that: it is outside the superellipse family, not its $p\to\infty$
limit. The single-angle ratio cannot distinguish the box from beyond-the-box;
only the full contour can.

\end{document}
